% =====================================================================
%  Chapter 19 review -- the 14-station revision contract.
%  Revision-only: tables, one-liners, pointers back via \secref/\thmref/
%  \propref/\defref/\coderef/\figref/\probref.  No re-derivations.
%  Summit: CF 414D Mashmokh and Water Tanks -- a litre at depth d reaches the
%  root d moves after release; to win X dollars land X litres on ONE move.
%  Binary-search the largest X <= min(k, m-1) whose synchronisation cost
%      min over level r of  ( X*r - sum of the X deepest depths among tanks
%                             at depth <= r )  <= p
%  is affordable; the inner "X deepest depths up to r" is an order-statistic
%  DESCEND (19.3, lst:c13-kthseg's shape) over a depth-axis tree storing the
%  (count, sum-of-depths) record (19.2, thm:c13-subarray-style tuple node).
%  Samples verified live against the CF statement:
%    (1) m=10 k=2 p=1  (edges)              -> 2   (X capped by k=2)
%    (2) m=5  k=1000 p=1000 (edges)         -> 4   (all 4 non-root tanks)
%  The official editorial frames it as a linear two-pointer over depth levels;
%  we give the segment-tree reading and note that alternative honestly.
%  Exports prob:c13-summit, lst:c13-kit, lst:c13-summit.
% =====================================================================

\chaptersection{Chapter 19 Summary}

One sentence: \emph{a segment tree is a binary tree whose nodes are intervals of
an index axis and whose payload is any associative fold with an identity, so a
single \(\bigO{\log n}\) walk folds a range and a second walks a point update to
the root -- and every question Chapter~18's Fenwick tree had to refuse because its
combine was not an invertible sum (range \(\min/\max/\gcd\) under arbitrary writes,
a maximum subarray, "the first position clearing a threshold") is one canonical
cover or one root-to-leaf descent away.} Chapter~18 kept an array that hid a tree
over the \emph{lowbit} blocks; this chapter keeps the whole tree and pays a modest
constant for the power to combine \emph{anything} associative. Each row is a
pointer, not a proof.

\begin{center}
\renewcommand{\arraystretch}{1.25}
\begin{tabular}{@{}p{0.24\linewidth} p{0.48\linewidth} p{0.20\linewidth}@{}}
\toprule
\textbf{Result} & \textbf{Statement} & \textbf{Where} \\
\midrule
Node and interval & Each node \(v\) is tagged with a closed interval \([tl,tr]\);
its children split at the midpoint, in a \(4n\) heap. & \defref{c13-node} \\
Merge contract & An \emph{associative} fold \(\mathcal{M}\) with identity \(e\)
makes build/update/query correct and \(\bigO{\log n}\). & \thmref{c13-merge} \\
Identity element & A fold needs \(e\): \(\min\!\to\!+\infty\), \(\max\!\to\!-\infty\),
sum/gcd/xor\(\to0\); a node may be a whole \emph{record}. & \defref{c13-identity} \\
Four-field merge & Max-subarray stores \((\mathit{tot},\mathit{pref},\mathit{suff},
\mathit{best})\); merge is \(O(1)\) but \emph{not} commutative. & \thmref{c13-subarray} \\
Monotone descend & If an aggregate decides "is there a target below?", walk
root-to-leaf, pruning whole subtrees, in \(\bigO{\log n}\). & \thmref{c13-descend} \\
\(2n\) iterative & Complete tree in a flat array: one build loop, one climbing
update, one two-accumulator half-open query. & \propref{c13-iter} \\
Static + offline & Build once, then every read-only fold or threshold descend is
an \(\bigO{\log n}\) oracle; sort the queries to reuse one pass. & \propref{c13-offline} \\
\bottomrule
\end{tabular}
\end{center}

\medskip
{\small\textbf{Algorithms from this chapter.}}
\begin{center}
\renewcommand{\arraystretch}{1.2}
\begin{tabular}{@{}p{0.27\linewidth} l l >{\raggedright\arraybackslash}p{0.26\linewidth}@{}}
\toprule
\textbf{Algorithm} & \textbf{Prep} & \textbf{Per op} & \textbf{Code} \\
\midrule
Build / point update / range fold & \(\bigO{n}\) build & \(\bigO{\log n}\) & \coderef{c13-build}, \coderef{c13-update}, \coderef{c13-query} \\
CSES 1649 (Dynamic Range Min) & \(\bigO{n}\) & \(\log n\) & \coderef{c13-1649-soln} \\
Max-subarray tuple merge & \(\bigO{n}\) & \(\log n\) & \coderef{c13-combine} \\
CSES 1190 (Subarray Sum Queries) & \(\bigO{n}\) & \(\log n\) & \coderef{c13-1190-soln} \\
gcd / count-of-max combine & \(\bigO{n}\) & \(\log n\) & \coderef{c13-gcdcount} \\
Descend: first prefix \(\ge\) target & \(\bigO{n}\) & \(\log n\) & \coderef{c13-descend} \\
\(k\)-th live item by counts (descend) & \(\bigO{n}\) & \(\log n\) & \coderef{c13-kthseg} \\
CF 1179C (Serge and Dining Room) & \(\bigO{n}\) & \(\log n\) & \coderef{c13-1179-soln} \\
Iterative \(2n\) tree (build/update/query) & \(\bigO{n}\) & \(\log n\) & \coderef{c13-iter}, \coderef{c13-iter-range} \\
CF 514D (R2D2 and Droid Army) & \(m\) trees & \(\log n\) & \coderef{c13-514-soln} \\
Static leftmost \(\ge x\) (no updates) & \(\bigO{n}\) & \(\log n\) & \coderef{c13-firstge} \\
CF 940E (Cashback, DP + range-min) & \(\bigO{n}\) & \(\log n\) & \coderef{c13-940-soln} \\
CF 414D (Summit, Water Tanks) & \(\bigO{n}\) & \(\log n\) & \coderef{c13-summit} \\
\bottomrule
\end{tabular}
\end{center}

\chaptersection{Master Decision Flowchart}

Post-study ladder: walk top to bottom, stop at the first branch that matches. It
decides \emph{whether a segment tree is even the right object} and, if so,
\emph{which flavour}, then hands off what this chapter does \emph{not} own: range
updates with a lazy tag (Ch.~20 \(\chref{c15}\)), \(k\)-th-value-in-a-static-range
and persistence (Ch.~21 \(\chref{c19}\)), and path queries on a tree
(Ch.~22 \(\chref{c16}\)).

\begin{lstlisting}[language={}, basicstyle=\ttfamily\small, numbers=none, frame=single]
  You need a live RANGE aggregate over an index axis.  Which structure?
  |
  |-- is the combine a plain SUM (invertible) with point adds?
  |     +-- no segtree needed: a prefix array (Ch.5) or a BIT (Ch.18, c14)
  |        is smaller and answers a range sum as two prefixes.
  |
  |-- non-invertible fold -- min / max / gcd / count -- with a point SET?
  |     +-- recursive segtree (19.1, lst:c13-build / -update / -query).
  |        CSES 1649 range-min-with-set is the hello-world (prob:c13-dynmin);
  |        this is the fold a BIT structurally cannot maintain.
  |
  |-- the answer is a RECORD, not a scalar (max subarray, brackets, pair-best)?
  |     +-- tuple node + O(1) ASSOCIATIVE merge (19.2, lst:c13-combine).
  |        Keep the merge order: it is NOT commutative (thm:c13-subarray).
  |
  |-- "first / last position whose aggregate crosses a threshold x"?
  |     +-- DESCEND root-to-leaf, pruning on the stored aggregate
  |        (19.3, lst:c13-descend; 19.5 static form lst:c13-firstge).
  |        A fold returns a value; a descend returns an index -- a BIT cannot
  |        see inside a range at all.
  |
  |-- "k-th live item / select by counts"?
  |     +-- count segtree + descend (19.3, lst:c13-kthseg) -- same answer as
  |        Ch.18's lowbit-lift, but by walking the tree, not lifting bits.
  |
  |-- fixed length, commutative fold, want the smallest constant and least code?
  |     +-- iterative 2n tree (19.4, lst:c13-iter + lst:c13-iter-range):
  |        half-open [l,r), two accumulators if the fold is order-sensitive.
  |
  |-- array NEVER changes; many read-only folds / descends; a DP transition?
  |     +-- build ONCE, query as an oracle (19.5, prop:c13-offline).
  |        Pure static min/max with ZERO updates?  a sparse table (Ch.21) is
  |        O(1) per query -- but you then lose updates, descends and records.
  |
  |-- a WHOLE INTERVAL is modified (add / set a range), not just a point?
  |        +-- NOT this chapter: lazy propagation, Ch.20 (c15).
  |-- "k-th smallest in a STATIC range", or persistent versions?
  |     +-- merge-sort tree / persistent segtree, Ch.21 (c19).
  |-- path queries on a TREE?
        +-- flatten with HLD onto a segtree, Ch.22 (c16).
\end{lstlisting}

\chaptersection{Common Mistakes Checklist}

\begin{itemize}[nosep, leftmargin=1.4em]
  \item \textbf{Under-sizing the array.} The \(4n\) heap layout of
        \(\defref{c13-node}\) needs \(4n\), not \(2n\), slots for the \emph{recursive}
        tree whenever \(n\) is not a power of two; the \(2n\) form of
        \(\coderef{c13-iter}\) is legal \emph{only} because its tree is complete.
        A recursive tree in a \(2n\) array runs off the end (RE).
  \item \textbf{Wrong or missing identity element.} A range query folds the
        \emph{disjoint} branches to \(e\): \(\min\) needs \(+\infty\), \(\max\) needs
        \(-\infty\), sum/gcd/xor need \(0\). Returning a real leaf value (\(\coderef{c13-query}\))
        on a disjoint interval silently corrupts the fold (\(\defref{c13-identity}\)).
  \item \textbf{Half-open vs closed range drift.} \(\coderef{c13-query}\) folds a
        \emph{closed} \([L,R]\); \(\coderef{c13-iter-range}\) folds a \emph{half-open}
        \([l,r)\). Mixing the two conventions mis-answers by one at every boundary.
  \item \textbf{Assuming the combine is commutative.} The max-subarray merge of
        \(\thmref{c13-subarray}\) and the bracket triple are associative but \emph{not}
        commutative, so the iterative query must keep the two accumulators in
        left/right order; a single accumulator passes \(\min\) and fails those.
  \item \textbf{Descend going the wrong way.} \(\coderef{c13-descend}\)'s and
        \(\coderef{c13-firstge}\)'s answer is the \emph{leftmost} qualifier: probe the
        left (or right, for a "largest") child \emph{first} and only fall through when
        it fails. Probing the wrong side returns a valid-but-not-extremal index.
  \item \textbf{\(2n\) padding.} Skipping the power-of-two pad is fine for a
        commutative fold; for an order-sensitive one pad to \(2^{\lceil\log_2 n\rceil}\)
        with \(e\) so every leaf sits at the same depth (\(\propref{c13-iter}\)).
  \item \textbf{64-bit sums.} A sum tree over \(n=2\cdot10^5\) values \(\le10^9\)
        overflows \code{int} at the root; the payload array is \code{long long} the
        moment the fold can exceed \(2^{31}\) (\(\coderef{c13-build}\)).
  \item \textbf{Building with \(n\) updates.} The bottom-up fold
        (\(\coderef{c13-build}\), \(\coderef{c13-iter}\)) is \(\bigO{n}\); calling
        \(\coderef{c13-update}\) per element is \(\bigO{n\log n}\) and the TLE when the
        tree sits inside an outer binary search (\(\probref{c13-r2d2}\)).
  \item \textbf{Reaching for a segtree on a plain sum.} If the fold is invertible,
        Ch.~18's BIT or Ch.~5's prefix array already solves it with less code and a
        smaller constant -- that is the boundary this chapter exists to cross, not
        re-litigate.
\end{itemize}


\chaptersection{WA Clinic}

Verdicts are symptoms; the causes are the short list this chapter drilled.
Diagnose top-down -- every leaf names a listing.

\begin{figure}[H]
\centering
\begin{tikzpicture}[diagram,
    v/.style={boxnode, fill=primaryblue!12, draw=primaryblue, text width=22mm, align=center},
    b/.style={boxnode, fill=difforange!12, draw=difforange, text width=12mm},
    l/.style={boxnode, fill=successgreen!10, draw=successgreen, text width=48mm,
              font=\scriptsize, align=left}]
  \node[v] (verdict) {\textbf{Verdict on the judge}};
  \node[b, below left=8mm and 5mm of verdict] (wa) {\textbf{WA}};
  \node[b, below right=8mm and 5mm of verdict] (tle) {\textbf{TLE}};
  \node[l, below right=1mm and 2mm of wa.south] (wa1)
    {disjoint branch returns a leaf value, not the identity \(e\): a \(\min\) fold
     answers \(0\), not \(+\infty\) (\(\coderef{c13-query}\))};
  \node[l, below left=1mm and 2mm of wa.south] (wa2)
    {closed \([L,R]\) coded as half-open, or the descend probes the wrong side
     (\(\coderef{c13-descend}\), \(\coderef{c13-iter-range}\))};
  \node[l, below right=1mm and 2mm of tle.south] (tle1)
    {recursive tree sized \(2n\) not \(4n\): a root path runs off the
     array (\(\defref{c13-node}\))};
  \node[l, below left=1mm and 2mm of tle.south] (tle2)
    {\(\bigO{n\log n}\) build inside a binary search; \code{int} sum overflow
     (\(\coderef{c13-iter}\))};
  \draw[flow] (verdict) -- (wa);
  \draw[flow] (verdict) -- (tle);
  \draw[flow] (wa) -- (wa1);
  \draw[flow] (wa) -- (wa2);
  \draw[flow] (tle) -- (tle1);
  \draw[flow] (tle) -- (tle2);
\end{tikzpicture}
\caption*{WA: a non-identity on the empty side, or a range or descend
convention off by one. TLE or RE: an under-sized \(4n\) array, a build repeated
inside a search, or an \code{int} overflow.}
\end{figure}

\chaptersection{Quick Reference Card}

\begin{quickref}[Chapter 19 at a glance]
\textbf{The object.} A binary tree over closed intervals of an index axis
(\(\defref{c13-node}\)); each node stores a fold of its interval under an
associative \(\mathcal{M}\) with identity \(e\). Build \(\bigO{n}\), point update
and range query \(\bigO{\log n}\) (\(\thmref{c13-merge}\),
\(\coderef{c13-build}\)/\(\coderef{c13-update}\)/\(\coderef{c13-query}\)).

\medskip
\textbf{The fold.} \(\min\!\to\!e{=}{+}\infty\), \(\max\!\to\!{-}\infty\),
sum/xor/gcd\(\to0\). A node may be a \emph{record}: max-subarray is
\((\mathit{tot},\mathit{pref},\mathit{suff},\mathit{best})\) merged in \(O(1)\),
associative but \textbf{not} commutative (\(\defref{c13-identity}\),
\(\thmref{c13-subarray}\), \(\coderef{c13-combine}\)).

\medskip
\textbf{The search.} If a stored aggregate decides "no target in this subtree",
walk root-to-leaf and prune: first prefix \(\ge\) target, leftmost \(\ge x\) in a
range, \(k\)-th by counts (\(\thmref{c13-descend}\),
\(\coderef{c13-descend}\)/\(\coderef{c13-kthseg}\)/\(\coderef{c13-firstge}\)).
This is binary search whose array is the tree -- Ch.~18's lowbit-lift is the
cousin (\(\secref{c13-search}\)).

\medskip
\textbf{\(2n\) form.} Complete tree in a flat array: one build loop, one climbing
update \code{p>>=1} with sibling \code{p\textasciicircum1}, one half-open query
\([l,r)\) using two order-keeping accumulators (\(\propref{c13-iter}\),
\(\coderef{c13-iter}\)/\(\coderef{c13-iter-range}\)).

\medskip
\textbf{Static/offline.} Build once, answer read-only folds/descends as an
\(\bigO{\log n}\) oracle; sort the queries and sweep to reuse one pass
(\(\propref{c13-offline}\)). Pure static \(\min/\max\) with no updates: a sparse
table is \(O(1)\) but gives up updates, descends, and records (Ch.~21 \(\chref{c19}\)).

\medskip
\textbf{Translation.} "range min/max/gcd under point set" \(\to\) recursive tree.
"max subarray / brackets / pair-best" \(\to\) tuple node + ordered merge. "first /
leftmost / \(k\)-th" \(\to\) descend. "fixed length, tight constant" \(\to\) \(2n\)
tree. "array never changes, many lookups" \(\to\) build once, query as oracle.
"add to a whole range" \(\to\) \emph{not here}: lazy, Ch.~20 (\(\chref{c15}\)).
\end{quickref}

\begin{lstlisting}[caption={Chapter 19 kit: the recursive tree, the iterative 2n form, and the max-subarray combine, complete and production-shaped}, label={lst:c13-kit}]
// ---- recursive segment tree (4n): ANY associative fold M with identity e ----
struct Seg{
    int n; vector<long long> t; long long (*M)(long long,long long); long long e;
    Seg(int n, long long (*M)(long long,long long), long long e)
        : n(n), t(4*n+4, e), M(M), e(e) {}
    void build(vector<long long>&a,int v,int l,int r){              // O(n)
        if(l==r){ t[v]=a[l]; return; }
        int m=(l+r)/2; build(a,2*v,l,m); build(a,2*v+1,m+1,r);
        t[v]=M(t[2*v],t[2*v+1]);
    }
    void update(int pos,long long val,int v,int l,int r){           // O(log n) point SET
        if(l==r){ t[v]=val; return; }
        int m=(l+r)/2;
        if(pos<=m) update(pos,val,2*v,l,m); else update(pos,val,2*v+1,m+1,r);
        t[v]=M(t[2*v],t[2*v+1]);
    }
    long long query(int L,int R,int v,int l,int r){                 // CLOSED [L,R] fold
        if(L>r || R<l) return e;                                    // disjoint -> identity
        if(L<=l && r<=R) return t[v];                               // covered -> whole node
        int m=(l+r)/2;
        return M(query(L,R,2*v,l,m), query(L,R,2*v+1,m+1,r));
    }
};
// ---- iterative 2n tree: fixed length, commutative (or two-accumulator) fold ----
int N; vector<long long> bi;                                        // t[1..2N-1], leaves t[N..2N-1]
long long mergeOp(long long a,long long b){ return min(a,b); }      // swap the fold here
const long long ID=(1LL<<60);                                       // identity e for THAT fold
void ibuild(vector<long long>&a){ N=(int)a.size(); bi.assign(2*N,ID);
    for(int i=0;i<N;i++) bi[N+i]=a[i];
    for(int i=N-1;i>0;i--) bi[i]=mergeOp(bi[2*i],bi[2*i+1]); }
void iupd(int p,long long val){ for(bi[p+=N]=val; p>1; p>>=1)
    bi[p>>1]=mergeOp(bi[p], bi[p^1]); }                             // p^1 = sibling
long long iqry(int l,int r){ long long L=ID,R=ID;                   // HALF-OPEN [l,r)
    for(l+=N,r+=N; l<r; l>>=1, r>>=1){
        if(l&1) L=mergeOp(L,bi[l++]);                               // left edge, L->R order
        if(r&1) R=mergeOp(bi[--r],R);                               // right edge, kept ordered
    } return mergeOp(L,R); }
// ---- max-subarray node: the 4-field combine a BIT cannot maintain ----
struct Node{ long long tot,pref,suff,best; };
Node mergeNode(Node a,Node b){ Node r;                              // associative, NOT commutative
    r.tot  = a.tot + b.tot;
    r.pref = max(a.pref, a.tot + b.pref);
    r.suff = max(b.suff, b.tot + a.suff);
    r.best = max( max(a.best, b.best), a.suff + b.pref );
    return r; }
// identity Node: tot=0, pref=suff=best = a very small number (-inf)
\end{lstlisting}


\chaptersection{Write-From-Memory Drills}

Five signature skeletons, each lifted from a listing already in the chapter.
Cover the marked line, reproduce from memory, then check the answer one line
below.

\begin{lstlisting}[caption={Drill 1 -- build: fold leaves up to the root (from lst:c13-build)}]
void build(vector<long long>& a, int v, int l, int r) {
    if (l == r) { t[v] = a[l]; return; }         // leaf: the singleton interval
    int m = (l + r) / 2;
    build(a, 2*v, l, m); build(a, 2*v+1, m+1, r);
    t[v] = M(t[2*v], t[2*v+1]);                  // one merge per internal node -> O(n)
}
\end{lstlisting}
\noindent{\scriptsize Key: every internal node is exactly the fold of its two
children, so a single post-order pass writes all \(2n\) nodes in \(\bigO{n}\) --
never build by calling \(\coderef{c13-update}\) \(n\) times (\(\thmref{c13-merge}\)).}

\begin{lstlisting}[caption={Drill 2 -- point update: set A[pos]=val, retrace one root path (from lst:c13-update)}]
void update(int pos, long long val, int v, int l, int r) {
    if (l == r) { t[v] = val; return; }          // SET, not add: a segtree does writes
    int m = (l + r) / 2;
    if (pos <= m) update(pos, val, 2*v,   l, m);
    else          update(pos, val, 2*v+1, m+1, r);
    t[v] = M(t[2*v], t[2*v+1]);                  // re-fold on the way back up
}
\end{lstlisting}
\noindent{\scriptsize Key: only the \(\bigO{\log n}\) ancestors of the leaf change;
a \emph{set} (not a delta) is precisely what a BIT cannot do for \(\min/\max\)
(\(\probref{c13-dynmin}\), \(\coderef{c13-update}\)).}

\begin{lstlisting}[caption={Drill 3 -- closed-range query: fold the canonical cover (from lst:c13-query)}]
long long query(int L, int R, int v, int l, int r) {   // fold over CLOSED [L,R]
    if (L > r || R < l) return e;                        // disjoint -> identity
    if (L <= l && r <= R) return t[v];                   // covered -> take the node
    int m = (l + r) / 2;
    return M(query(L,R,2*v,l,m), query(L,R,2*v+1,m+1,r));
}
\end{lstlisting}
\noindent{\scriptsize Key: the two base cases are the whole method -- a disjoint
branch contributes \(e\), a covered branch is folded as one node; the recursion
visits \(\bigO{\log n}\) nodes (\(\thmref{c13-merge}\), \(\defref{c13-identity}\)).}

\begin{lstlisting}[caption={Drill 4 -- max-subarray merge: the four-field combine (from lst:c13-combine)}]
Node mergeNode(Node a, Node b) {                 // LEFT child a, RIGHT child b
    Node r;
    r.tot  = a.tot + b.tot;
    r.pref = max(a.pref, a.tot + b.pref);        // best prefix stays left or crosses
    r.suff = max(b.suff, b.tot + a.suff);        // best suffix stays right or crosses
    r.best = max( max(a.best, b.best), a.suff + b.pref );  // or spans the seam
    return r;
}
\end{lstlisting}
\noindent{\scriptsize Key: the merge is \textbf{order-sensitive} -- \(a\) must be
the left interval -- and the identity node has \(\mathit{tot}{=}0\), the rest
\(-\infty\); this combine has no BIT analogue (\(\thmref{c13-subarray}\)).}

\begin{lstlisting}[caption={Drill 5 -- descend to the first prefix >= target (from lst:c13-descend)}]
int descend(long long target, int v, int l, int r) {   // SUM tree; find least p
    if (l == r) return l;                               // with prefix[1..p] >= target
    int m = (l + r) / 2;
    if (t[2*v] >= target) return descend(target, 2*v, l, m);        // answer is left
    return descend(target - t[2*v], 2*v+1, m+1, r);                 // else right, minus left
}
\end{lstlisting}
\noindent{\scriptsize Key: the stored aggregate decides which half can hold the
target, so one child is pruned per level -- \(\bigO{\log n}\) descent, the tree's
form of binary search (\(\thmref{c13-descend}\)); this is the cousin of Ch.~18's
lowbit-lift (\(\secref{c13-search}\)).}

\chaptersection{Complete Problem Checklist}

Every problem this chapter offers -- 4 fully solved in-section, worked entries,
practice, and one Summit. Tick \emph{only} after an unaided solve; a hinted
attempt counts as half.

\begin{center}
\footnotesize
\renewcommand{\arraystretch}{1.25}
\begin{tabular}{@{}c p{0.19\linewidth} l p{0.20\linewidth} p{0.24\linewidth} c@{}}
\toprule
\textbf{\#} & \textbf{Problem} & \textbf{Source} & \textbf{Key technique} & \textbf{What to look for} & \textbf{Done?} \\
\midrule
\multicolumn{6}{l}{\textbf{Section 19.1 -- the segment tree core}} \\
W19.1 & Dynamic Range Minimum Queries & CSES 1649 & \lv{3} min-tree, point set & set \(=\) write to leaf, retrace one path; disjoint \(=+\infty\) & \(\square\) \\
\midrule
\multicolumn{6}{l}{\textbf{Section 19.2 -- arbitrary combine}} \\
W19.2 & Subarray Sum Queries & CSES 1190 & \lv{4} four-field tuple node & \((\mathit{tot},\mathit{pref},\mathit{suff},\mathit{best})\); non-commutative merge & \(\square\) \\
W19.2b & gcd / count-of-max combine & worked & \lv{4} swap the fold & same skeleton, new \(\mathcal{M}\); identity for gcd is \(0\) & \(\square\) \\
\midrule
\multicolumn{6}{l}{\textbf{Section 19.3 -- searching the tree}} \\
W19.3 & Serge and Dining Room & CF 1179C & \lv{5} suffix-max tree + descend & descend to first prefix \(\le x\); two point writes per move & \(\square\) \\
W19.3b & \(k\)-th live item by counts & worked & \lv{4} count-tree descend & prune on subtree count; contrast Ch.18 lift & \(\square\) \\
\midrule
\multicolumn{6}{l}{\textbf{Section 19.4 -- iterative tree}} \\
W19.4 & R2D2 and Droid Army & CF 514D & \lv{5} \(m\) max-trees + binary search & half-open \([l,r)\); monotone feasibility; print the \(m\) maxima & \(\square\) \\
\midrule
\multicolumn{6}{l}{\textbf{Section 19.5 -- static / offline}} \\
W19.5a & leftmost \(\ge x\) in a range & worked & \lv{4} static max-tree descend & build once; prune on \(\mathit{mx}<x\); left child first & \(\square\) \\
W19.5b & Cashback & CF 940E & \lv{5} range-min oracle + DP & optimal runs length \(1\) or \(c\); window min is a static query & \(\square\) \\
\midrule
\multicolumn{6}{l}{\textbf{Summit (chapter end)}} \\
W19 & Mashmokh and Water Tanks & CF 414D & \lv{5} binary-search answer + order-statistic descend & depth-axis record tree; cost \(Xr - \sum\) deepest \(X\) & \(\square\) \\
\bottomrule
\end{tabular}
\end{center}

\chaptersection{Recommended Study Order}

Five focused days. Day~1 internalises the build/update/query triple; day~2 the
tuple record and its ordered merge; day~3 the descend and its Ch.~18 contrast;
day~4 the iterative \(2n\) form and the BIT-vs-segtree accounting; day~5 the
static/offline oracle, the DP pairing, and the Summit.

\begin{center}
\small
\renewcommand{\arraystretch}{1.25}
\begin{tabular}{@{}llp{0.55\linewidth}@{}}
\toprule
\textbf{Day} & \textbf{Focus} & \textbf{Tasks} \\
\midrule
1 & Read + do 19.1 & re-trace build, one point update, one range query on the shared array \(A\) (\figref{c13-tree}) by hand; write \coderef{c13-build}/\coderef{c13-update}/\coderef{c13-query} cold; CSES~1649 end to end \\
2 & Read + do 19.2 & derive the four-field merge from scratch; re-derive the identity node; CSES~1190; then swap the fold to gcd and count-of-max (\(\coderef{c13-gcdcount}\)) \\
3 & Read + do 19.3 & the first-prefix \(\ge\) descend on \(A\); the \(k\)-th-by-counts descend; side-by-side with Ch.~18's lowbit-lift; CF~1179C \\
4 & Read + do 19.4 & the \(2n\) build/update/query loops; the two-accumulator half-open walk; fill the BIT / recursive / iterative cost table; CF~514D \\
5 & Do 19.5 + Summit & static leftmost-\(\ge x\) oracle (\(\coderef{c13-firstge}\)); the 940E length-\(1\)-or-\(c\) DP; then CF~414D; walk the flowchart; first 3 mixed \\
\bottomrule
\end{tabular}
\end{center}


\chaptersection{Difficulty Progression}

Tiers are competency gates -- each row names the \emph{skill} the tier unlocks.

\begin{center}
\small
\renewcommand{\arraystretch}{1.3}
\begin{tabular}{@{}lp{0.34\linewidth}p{0.38\linewidth}@{}}
\toprule
\textbf{Tier} & \textbf{Problems} & \textbf{What passing this tier proves} \\
\midrule
\lv{3}--\lv{4} \ \(\bigstar\bigstar\bigstar\bigstar\)
  & CSES~1649, worked gcd/count swap, static leftmost-\(\ge x\)
  & you build the tree cold, drive the point update and range fold, and change the \emph{stored quantity} without touching the skeleton \\
\lv{4}--\lv{5} \ \(\bigstar\bigstar\bigstar\bigstar\bigstar\)
  & CSES~1190, CF~1179C, CF~514D, CF~940E
  & you store a \emph{record} per node, keep an ordered (non-commutative) merge straight, and turn a threshold question into a descend \\
\lv{5}{+} \ \(\bigstar\bigstar\bigstar\bigstar\bigstar\)
  & CF~414D (Summit), CF~1208E, CF~515E, CF~500E, CF~675E
  & you compose a segtree with an outer binary search or a DP, descend for an order statistic, and harden for \(n\le10^5\), values \(\le10^{15}\), \(k,p\le10^{9}\) \\
\bottomrule
\end{tabular}
\end{center}

\chaptersection{Reduction Gallery}

Every entry is ``one segtree flavour + one twist.'' Each line collapses a problem
family onto a listing already in the chapter.

\begin{itemize}[nosep, leftmargin=1.4em]
  \item ``range min/max with a point \emph{set}'' (CSES~1649) \(\to\) recursive
        tree, write to leaf, retrace one path (\(\coderef{c13-update}\),
        \(\probref{c13-dynmin}\)) -- the fold a BIT structurally refuses.
  \item ``best subarray under edits'' (CSES~1190) \(\to\) four-field record merged
        in \(O(1)\), order-sensitive (\(\thmref{c13-subarray}\), \(\coderef{c13-combine}\)).
  \item ``range gcd / how many equal the range max'' (worked) \(\to\) same skeleton,
        new associative \(\mathcal{M}\) (\(\coderef{c13-gcdcount}\), \(\defref{c13-identity}\)).
  \item ``first position with prefix \(\ge\) target / \(k\)-th live item'' (CF~1179C,
        hotel) \(\to\) descend the tree, prune on the stored aggregate
        (\(\thmref{c13-descend}\), \(\coderef{c13-descend}\)/\(\coderef{c13-kthseg}\)).
  \item ``longest window whose per-column maxima fit a budget'' (CF~514D) \(\to\)
        \(m\) iterative max-trees + a binary search on the length
        (\(\coderef{c13-iter}\), \(\probref{c13-r2d2}\)).
  \item ``partition so each run drops its smallest'' (CF~940E) \(\to\) DP whose
        transition is a static range-min oracle (\(\coderef{c13-940-soln}\),
        \(\probref{c13-cashback}\), \(\propref{c13-offline}\)).
  \item ``land X water on one move, cheapest'' (CF~414D) \(\to\) binary-search the
        answer + an order-statistic descend over a depth-axis record tree
        (\(\probref{c13-summit}\), \(\coderef{c13-summit}\)).
\end{itemize}

\chaptersection{Mixed Problem Set}

Ten entries from CF~1900 to CF~2300, all verified against the live problem set
and absent from the rest of the book. None says ``use a segment tree'' --
recognise the flavour from what the problem folds, descends to, or sweeps. Each
hint names its reduction and the chapter listing it leans on.

\begin{mixedset}
  \practice{Bash and a Tough Math Puzzle}{CF 914D}
    {https://codeforces.com/problemset/problem/914/D}{CF 1900 L3}
    {Type-1 query asks whether the \(\gcd\) of \([l,r]\) can be made a multiple of
     \(x\) by changing \emph{at most one} element; type-2 is a point set. A range
     \(\gcd\) fold (identity \(0\), \(\coderef{c13-gcdcount}\)) plus a
     \emph{descend} that counts elements not divisible by \(x\) and stops the moment
     it exceeds one -- prune a whole subtree when its stored \(\gcd\) is already
     divisible (\(\thmref{c13-descend}\), \(\coderef{c13-descend}\)) (\(\secref{c13-combine}\), \(\secref{c13-search}\)).}
  \practice{Pillars}{CF 474E}
    {https://codeforces.com/problemset/problem/474/E}{CF 2000 L4}
    {Longest jump sequence where consecutive pillars differ in height by at least
     \(d\): an LIS-style \(\mathit{dp}[j]=1+\max\{\mathit{dp}[i]:h_i\le h_j-d\) or
     \(h_i\ge h_j{+}d\}\). Compress heights to ranks and keep a \emph{max}-segtree over
     the height axis: each transition is a range-max fold + a point update
     (\(\coderef{c13-query}\), \(\coderef{c13-update}\)); store parent pointers to
     rebuild the sequence (\(\secref{c13-core}\)).}
  \practice{Babaei and Birthday Cake}{CF 629D}
    {https://codeforces.com/problemset/problem/629/D}{CF 2000 L4}
    {Stack cylinders in index order, each strictly larger in volume than the one
     below, maximising total volume: \(\mathit{dp}[i]=\mathit{vol}[i]+\max\{\mathit{dp}[j]:j<i,\ \mathit{vol}[j]<\mathit{vol}[i]\}\).
     Compress volumes, drive a prefix range-max + point update on a segtree
     (\(\coderef{c13-build}\), \(\coderef{c13-query}\)) -- the same max-fold as
     Pillars, one axis per coordinate (\(\secref{c13-core}\)).}
  \practice{SUM and REPLACE}{CF 920F}
    {https://codeforces.com/problemset/problem/920/F}{CF 2000 L4}
    {Range update \(a_i\leftarrow D(a_i)\) (divisor count) and range-sum query. No
     lazy tag is needed: store \((\text{sum},\max)\) per node and \emph{descend} only
     into leaves whose value is not yet a fixed point (\(\max\le2\) prunes the whole
     subtree, since \(D(1){=}1,D(2){=}2\)) -- a \(\secref{c13-search}\) descend on the
     \(\max\) aggregate, sum fold from \(\coderef{c13-combine}\).}
  \practice{Gifts Order}{CF 2057D}
    {https://codeforces.com/problemset/problem/2057/D}{CF 2000 L4}
    {Maximise \(\max(a_l..a_r)-\min(a_l..a_r)-(r-l)\) over all intervals after each
     point set. Rewrite the span term as a per-node \emph{record}: keep
     \(\max/\min\) of \(a_i-i\) and of \(a_i+i\) and merge in \(O(1)\) exactly like
     the four-field max-subarray node (\(\thmref{c13-subarray}\),
     \(\coderef{c13-combine}\)); point update retraces one path (\(\coderef{c13-update}\)).}
  \practice{Friends and Subsequences}{CF 689D}
    {https://codeforces.com/problemset/problem/689/D}{CF 2100 L4}
    {Count pairs \((l,r)\) with \(\max(a_l..a_r)=\min(b_l..b_r)\) over two fixed
     arrays. For a \emph{fixed} \(l\), \(\max a\) is non-decreasing and \(\min b\)
     non-increasing in \(r\), so their difference is monotone -- binary-search the
     \(r\)-window where they are equal, reading each \(\max/\min\) from a tree
     built \emph{once} on \(a\) and \(b\): the static range-fold oracle of
     \(\coderef{c13-build}\)/\(\coderef{c13-query}\) and \(\propref{c13-offline}\).
     A sparse table is the pure-static rival this section flags (\(\secref{c13-static}\)).}
  \practice{Let Them Slide}{CF 1208E}
    {https://codeforces.com/problemset/problem/1208/E}{CF 2200 L4}
    {Each row's array may slide within \(w\) cells; per column take the max
     contribution. For every row you need \(\max\) over a \emph{window of positions}
     that can reach column \(j\) -- a static range-max oracle (\(\coderef{c13-firstge}\)'
     build-once idea), accumulated onto the columns with a difference array; the whole
     thing is the ``array never changes, many read-only folds'' regime of
     \(\propref{c13-offline}\) (\(\secref{c13-static}\)).}
  \practice{Drazil and Park}{CF 515E}
    {https://codeforces.com/problemset/problem/515/E}{CF 2300 L5}
    {On a circular park, per query exclude one arc and maximise
     \(2(h_x+h_y)+\mathit{dist}(x,y)\) over a pair on the free arc. Double the array
     and give each node a \emph{tuple} (best pair, best prefix/suffix of the two
     derived quantities \(2h_i\pm\mathit{pref}_i\)) merged \(O(1)\) in order -- a
     non-commutative combine like \(\thmref{c13-subarray}\), queried over the
     wrap-around range (\(\coderef{c13-combine}\), \(\secref{c13-combine}\)).}
  \practice{Trains and Statistic}{CF 675E}
    {https://codeforces.com/problemset/problem/675/E}{CF 2300 L5}
    {Sum, over all station pairs, the minimum tickets to travel: a left-to-right
     greedy that needs, for each \(i\), the index \(k\le i\) with the \emph{maximum}
     reach \(a_k\). A static range-max-with-argmax oracle answers each step
     (\(\coderef{c13-firstge}\)' descend returns the \emph{position}, not just the
     value; \(\thmref{c13-descend}\)) folded into a running \(\mathit{dp}\) sum
     (\(\secref{c13-search}\), \(\secref{c13-static}\)).}
  \practice{New Year Domino}{CF 500E}
    {https://codeforces.com/problemset/problem/500/E}{CF 2300 L5}
    {Pushing domino \(x\) must reach \(y\); pay to lengthen. Precompute
     \(\mathit{nxt}[i]\) = the first domino \emph{not} covered by the chain from \(i\),
     which is a ``first position whose \(\max_j(p_j+l_j)\) clears a reach threshold''
     -- the descend of \(\coderef{c13-descend}\) on a max-segtree over the
     right-reach axis, then sum the gap costs along the jump chain
     (\(\thmref{c13-descend}\), \(\secref{c13-search}\)).}
\end{mixedset}


\chaptersection{The Summit -- CF 414D Mashmokh and Water Tanks}

\begin{notebox}[Why this is the chapter top rope -- and an honest caveat]
Sections 19.2 and 19.3 each own one half: 19.2 lets a node be a \emph{record}
folded by an associative merge; 19.3 lets you \emph{descend} to an order statistic
the tree can see but a BIT cannot. The Summit is both at once: it searches an
answer with a binary search and, inside every predicate, descends a depth-axis tree
of \((\text{count},\text{sum of depths})\) records for the \(X\) deepest usable
tanks. \textbf{Honest note:} the official editorial solves it in \(O(m)\) with a
\emph{two-pointer} over depth levels -- that is the intended route and it is faster.
This chapter asked for a segment-tree treatment, so we give one and pay a
\(\log\) for it; read the two-pointer as the optimisation exercise at the end.
\end{notebox}

\solvedproblem{c13-summit}{Mashmokh and Water Tanks}{CF 414D}
  {https://codeforces.com/problemset/problem/414/D}{CF 2300 L5}
  {binary-search the answer X + an order-statistic DESCEND over a depth-axis tree storing (count, sum-of-depths) records}

\begin{problemstatement}[CF 414D Mashmokh and Water Tanks]
A rooted tree of \(m\) vertices (vertex \(1\) is the root) is given; every vertex
holds a water tank, empty at the start. Mashmokh has \(k\) litres and \(p\) coins.
He first pours exactly one litre into each of \emph{at most} \(k\) chosen non-root
tanks. Then he performs moves until no tank holds water. In one move he may close
the door of any non-root tank holding \(w\) litres, paying \(w\) coins for that
move; a closed tank keeps its water, an open tank passes all its water up to its
parent. Tanks are processed shallowest-first (the root first, and the root is
\emph{emptied} at the start of every move), so a litre climbs at most one edge per
move. If \(l\) moves are made and \(w_i\) is the water in the root after move \(i\),
he wins \(\max(w_1,\dots,w_l)\) dollars. Report the maximum dollars he can guarantee.

\textbf{Input.} \(m,k,p\) (\(2\le m\le 10^{5}\); \(0\le k,p\le 10^{9}\)); then \(m-1\)
edges \(a_i,b_i\) (\(1\le a_i,b_i\le m\)).

\textbf{Output.} One integer -- the maximum dollars.

\textbf{Sample 1.} \(m{=}10,k{=}2,p{=}1\) with edges
\((1,2)\), \((1,3)\), \((3,4)\), \((3,5)\), \((2,6)\), \((6,8)\),
\((6,7)\), \((9,8)\), \((8,10)\) \(\to\) \texttt{2}.\\
\textbf{Sample 2.} \(m{=}5,k{=}1000,p{=}1000\) with edges
\((1,2),(1,3),(3,4),(3,5)\) \(\to\) \texttt{4}.
\end{problemstatement}

\solstep{Key Observations}
\begin{enumerate}[nosep,leftmargin=1.6em]
  \item Because the root is emptied at the start of every move, the win is the water
        that arrives \emph{on a single move}, never an accumulation. A litre released
        from depth \(d\) reaches the root exactly \(d\) moves later, so to win \(X\)
        dollars \(X\) litres must land on the \emph{same} move.
  \item Release a litre from depth \(d_i\) and it arrives at move \(d_i\). To make it
        arrive instead at move \(T\ge\max_i d_i\) he closes its (single-litre) tank for
        \(T-d_i\) early moves, costing \(T-d_i\) coins; synchronising \(X\) litres at
        level \(T\) therefore costs \(\sum_i(T-d_i)=XT-\sum_i d_i\).
  \item For a fixed number \(X\) and a fixed synchronising level \(r=T\), the cheapest
        choice is the \(X\) \emph{deepest} tanks at depth \(\le r\) (a deeper tank needs
        less holding). \(\text{Cost}(X)=\min_{r}\big(Xr-\!\sum(\text{the \(X\) deepest depths}\le r)\big)\).
  \item Feasibility ``\(\text{Cost}(X)\le p\) and \(X\le k\) and \(X\le m-1\)'' is
        \textbf{monotone} in \(X\) (each extra litre only raises the minimum cost), so
        the answer is the largest affordable \(X\) -- a binary search -- and each
        predicate needs an order-statistic \emph{descend}.
\end{enumerate}

\solstep{Reduction to a Descend over a Depth Axis}
Root the tree and bucket the \(m-1\) non-root vertices by depth; let
\(\mathit{cnt}[d]\) be the count at level \(d\) and \(D\) the maximum depth. Build a
segment tree on the closed axis \([0,D]\) whose leaf \(d\) stores the record
\((\mathit{cnt}[d],\,d\cdot\mathit{cnt}[d])\) and whose internal node is the
componentwise sum -- an associative, commutative merge with identity \((0,0)\),
exactly the record idea of \(\thmref{c13-subarray}\) and the \(O(n)\) build of
\(\coderef{c13-iter}\). A prefix array \(\mathit{cc}[r]=\sum_{d\le r}\mathit{cnt}[d]\)
answers ``how many tanks at depth \(\le r\)'' in \(O(1)\).

To evaluate \(\text{Cost}(X)\): for each level \(r\) with \(\mathit{cc}[r]\ge X\),
\emph{descend} the tree to collect the sum of the \(X\) largest depths that are
\(\le r\) -- visit the right (deeper) child first, swallow a whole node when its
\(\mathit{cnt}\le\) the remaining need, else recurse, stopping at the \(r\) cap
(\(\thmref{c13-descend}\), the sibling of \(\coderef{c13-kthseg}\)); the descend costs
\(\bigO{\log D}\). Take the minimum of \(Xr-\text{got}\) over \(r\). Binary-search \(X\)
over \([0,\min(k,m-1)]\) on the monotone predicate \(\text{Cost}(X)\le p\). Total
\(O(D\log m\log R)\) with \(D\le m\) and \(R\le m\) -- the editorial's two-pointer
collapses the outer two scans to a single \(O(m)\) pass.

\solstep{Trace on a Hand Instance}
\begin{tracebox}[Synchronise tanks of depth profile 1,2,3,3 with k=4 and p=2]
Non-root depths \(\{1,2,3,3\}\), so \(\mathit{cnt}[1]{=}1,\mathit{cnt}[2]{=}1,
\mathit{cnt}[3]{=}2\), \(D{=}3\), and \(\mathit{cc}=(0,1,2,4)\). The depth-axis tree
(root interval \([0,3]\)) stores \((\text{count},\text{sum})\) per node; the leaf
records are \((0,0),(1,1),(1,2),(2,6)\).

\textbf{Cost(4).} Only \(r{=}3\) has \(\mathit{cc}[3]{=}4\ge4\). Descend for the \(4\)
largest depths \(\le3\): the root count \(4\le4\) sits inside the cap, so swallow the
whole tree, \(\text{got}=1{+}2{+}3{+}3=9\). Cost \(=4\cdot3-9=3\). \(3>p{=}2\) \(\Rightarrow\) infeasible.

\textbf{Cost(3).} \(r{=}3\): descend \(3\) largest \(\le3\). The right child \([2,3]\)
holds count \(1{+}2=3\le3\), swallow it, \(\text{got}=2+6=8\), need satisfied. Cost
\(=3\cdot3-8=1\). (\(r{=}2\) has \(\mathit{cc}[2]{=}2<3\), skipped.) Min cost \(=1\le p\) \(\Rightarrow\) feasible.

So the largest affordable \(X\) is \(3\): hold the depth-\(2\) tank for one move so all
three land on move \(3\) -- keep tanks \(\{2,3,3\}\), closing the depth-\(2\) one once.
The binary search over \(X\in[0,4]\) returns \(3\).
\end{tracebox}

\noindent\textbf{Replaying the official samples.} \emph{Sample~1:} \(k=2\) caps
\(X\le2\); \(\text{Cost}(2)=0\le p=1\) (two level-\(2\) tanks land together on move~2,
needing no coins), so the answer is \(2\). \emph{Sample~2:} depths
\(\{1,2,3,3\}\) again with \(k=p=1000\), \(D=3\); all four non-root tanks are usable,
\(\text{Cost}(4)=4\cdot3-9=3\le1000\), and \(X=5\) is impossible (only \(m-1=4\) tanks),
so the answer is \(4\) -- matching both judge outputs.

\solstep{Implementation}
\begin{lstlisting}[caption={CF 414D: a depth-axis (count,sum) record tree, an order-statistic descend for the X deepest, wrapped in a binary search on the answer}, label={lst:c13-summit}]
#include <bits/stdc++.h>
using namespace std;
struct Node{ long long cnt=0, sum=0; };              // (tanks here, sum of their depths)
static inline Node mg(Node a,Node b){ return {a.cnt+b.cnt, a.sum+b.sum}; }  // assoc, commut
int H; long long CAP; vector<Node> t;
long long NEED, GOT;                                 // descend accumulators
void takeDeepest(int v,int l,int r){                 // grab NEED largest depths <= CAP
    if(NEED<=0 || l>CAP) return;                     // past the cap: nothing here
    if(r<=CAP && t[v].cnt<=NEED){ NEED-=t[v].cnt; GOT+=t[v].sum; return; }  // whole node
    if(l==r){ long long kk=min(NEED,t[v].cnt); NEED-=kk; GOT+=kk*l; return; }
    int m=(l+r)>>1;
    takeDeepest(2*v+1,m+1,r);                        // rightmost (deepest) child FIRST
    takeDeepest(2*v,l,m);
}
int main(){
    int m; long long k,p; scanf("%d%lld%lld",&m,&k,&p);
    vector<vector<int> > adj(m+1);
    for(int i=1,a,b;i<m;i++){ scanf("%d%d",&a,&b); adj[a].push_back(b); adj[b].push_back(a); }
    vector<int> dep(m+1,-1), bfs; bfs.push_back(1); dep[1]=0; int maxD=0;
    for(int s=0;s<(int)bfs.size();s++){ int u=bfs[s]; if(dep[u]>maxD) maxD=dep[u];
        for(int w:adj[u]) if(dep[w]<0){ dep[w]=dep[u]+1; bfs.push_back(w); } }   // BFS depths
    H=1; while(H<=maxD+1) H<<=1; t.assign(2*H, Node());
    for(int v=2;v<=m;v++){ t[H+dep[v]].cnt++; t[H+dep[v]].sum+=dep[v]; }
    for(int i=H-1;i>0;i--) t[i]=mg(t[i<<1], t[i<<1|1]);        // O(n) build (lst:c13-iter)
    vector<long long> cc(maxD+1,0);                            // #tanks at depth <= r
    for(int r=1;r<=maxD;r++) cc[r]=cc[r-1]+t[H+r].cnt;
    auto cost=[&](long long X)->long long{
        long long best=(long long)4e18;
        for(int r=1;r<=maxD;r++){
            if(cc[r]<X) continue;                               // not enough tanks <= r
            CAP=r; NEED=X; GOT=0; takeDeepest(1,0,H-1);         // sum of X deepest <= r
            long long c=X*r-GOT; if(c<best) best=c;             // hold cost to sync at level r
        }
        return best;
    };
    long long lo=0, hi=min(k,(long long)m-1), ans=0;            // monotone in X -> binary search
    while(lo<=hi){ long long mid=(lo+hi)/2;
        if(cost(mid)<=p){ ans=mid; lo=mid+1; } else hi=mid-1; }
    printf("%lld\n", ans);
}
// Depth-axis tree stores (count, sum-of-depths): the record-node merge of 19.2.  takeDeepest
// is the 19.3 order-statistic descend (right child first) that a BIT cannot run -- it reads
// the X largest depths among tanks at level <= CAP.  Binary-searching X on the monotone
// Cost(X) <= p is 19.4's outer-loop pattern.  X*r and GOT reach 1e10, so every accumulator is
// long long.  The official O(m) two-pointer is faster; this is the segment-tree reading.
\end{lstlisting}

\noindent\textbf{Complexity.} A BFS to bucket depths is \(\bigO{m}\); the \(2n\) build is
\(\bigO{D}\). Each predicate \(\text{Cost}(X)\) scans the \(D\) levels and runs one
\(\bigO{\log D}\) descend, so \(\bigO{D\log D}\); the answer binary search repeats it
\(\bigO{\log m}\) times for \(\bigO{D\log D\log m}\) overall -- fine inside \(m\le10^5\),
with the editorial's two-pointer the \(O(m)\) upgrade.

\begin{hardeningbox}[AC to red: hardening the Summit]
\begin{itemize}[nosep,leftmargin=1.4em]
  \item \textbf{The root is cleared every move.} The win is litres arriving on \emph{one}
        move, not the running total; do not model water as accumulating in the root.
  \item \textbf{Pick the \(X\) DEEPEST tanks, not shallowest.} Cost \(\sum(r-d_i)\) is
        minimised by large \(d_i\); taking the shallowest tanks inflates the hold cost
        and rejects feasible \(X\).
  \item \textbf{Guard \(\mathit{cc}[r]\ge X\) before descending.} With fewer than \(X\)
        tanks at depth \(\le r\) the descend returns a short sum and \(\text{Cost}(X)\)
        silently reads the wrong value -- the \code{continue} is load-bearing.
  \item \textbf{Monotonicity is what licenses the search.} \(\text{Cost}(X)\) is
        non-decreasing in \(X\); verify it before binary searching, or the standard
        \(\bigO{\log}\) answer trick breaks (\(\thmref{c13-descend}\)' predicate framing).
  \item \textbf{Levels start at \(1\).} Depth \(0\) is the root, which is never filled and
        never closed; scanning \(r\) from \(1\) skips it and keeps the cap honest.
  \item \textbf{64-bit throughout.} \(X\cdot r\) and the sum-of-depths reach
        \(10^{10}\); \code{int} overflows the cost and the sentinel \(\text{best}\).
\end{itemize}
\end{hardeningbox}

\begin{recallbox}[Chapter 19 in ten seconds]
A segment tree is a binary tree of \emph{intervals} whose node holds any associative
fold with an identity (\(\defref{c13-node}\), \(\thmref{c13-merge}\)): build
\(\bigO{n}\), point set and range fold \(\bigO{\log n}\) (\(\coderef{c13-build}\),
\(\coderef{c13-query}\)). Change the payload to a \emph{record} and you fold things a
BIT cannot -- the four-field max-subarray of \(\thmref{c13-subarray}\) is the type
example (\(\coderef{c13-combine}\)). Let the stored aggregate decide ``is the target
under me'' and you can \emph{descend} for the first/\(k\)-th qualifier in
\(\bigO{\log n}\) (\(\thmref{c13-descend}\), \(\coderef{c13-descend}\)) -- the tree's
answer to Ch.~18's lowbit-lift. Compact the recursion into a complete \(2n\) array and
you get the fastest, shortest version, at the price of a two-accumulator half-open walk
for order-sensitive folds (\(\propref{c13-iter}\), \(\coderef{c13-iter}\)). Build once
and never update, and the tree is a pure read-only oracle feeding offline sweeps and
DP transitions (\(\propref{c13-offline}\)).
\selftest{
  \item What two ingredients make a range query correct, and what breaks if the disjoint
        base case returns a leaf value instead of \(e\)? \hfill\lv{3}
  \item Why is the max-subarray merge associative but not commutative, and what does the
        iterative query do differently as a result? \hfill\lv{4}
  \item State the difference between a fold query and a descend query, with one problem
        each. \hfill\lv{4}
  \item In the Summit, which two chapter techniques compose, and what does each contribute? \hfill\lv{5}
}
\end{recallbox}


\chaptersection{Tool Kit Inventory}

The machine assembled through Chapter~19. Earlier chapters' rows are carried
forward; this chapter's are appended. Pointers only.

\begin{center}
\small
\renewcommand{\arraystretch}{1.25}
\begin{tabular}{@{}p{0.21\linewidth} l c p{0.25\linewidth} p{0.17\linewidth}@{}}
\toprule
\textbf{Tool} & \textbf{Since} & \textbf{Lv} & \textbf{Trigger phrase} & \textbf{Where} \\
\midrule
\multicolumn{5}{l}{\itshape Carried forward} \\
Euler tour \(\tin/\tout\) flatten & 5 & \lv{3} & ``a subtree is a contiguous range'' & \chref{c24} \\
Binary lifting / lowbit-lift select & 11, 18 & \lv{4} & ``the descend's bit-walking cousin'' & \chref{c11}, \coderef{c14-kth} \\
Invertible prefix sum & 5, 18 & \lv{2} & ``the fold that needs NO segtree'' & \chref{c14} \\
\midrule
\multicolumn{5}{l}{\itshape Chapter 19 -- new this chapter} \\
Recursive segtree (build/update/query) & 19.1 & \lv{3} & ``range min/max/gcd under point set'' & \coderef{c13-build}, \coderef{c13-query} \\
Arbitrary associative fold + identity & 19.2 & \lv{3} & ``a combine with no inverse'' & \thmref{c13-merge}, \defref{c13-identity} \\
Tuple/record node (max-subarray) & 19.2 & \lv{4} & ``best subarray / brackets / pair-best'' & \thmref{c13-subarray}, \coderef{c13-combine} \\
Descend to a threshold / \(k\)-th & 19.3 & \lv{4} & ``first position clearing a bound'' & \thmref{c13-descend}, \coderef{c13-descend} \\
Iterative \(2n\) tree (half-open) & 19.4 & \lv{3} & ``fixed length, tight constant, least code'' & \propref{c13-iter}, \coderef{c13-iter} \\
Static/offline range oracle & 19.5 & \lv{4} & ``array never changes, many lookups'' & \propref{c13-offline}, \coderef{c13-firstge} \\
Answer-search + record-descend & Summit & \lv{5} & ``binary-search the answer, descend inside'' & \coderef{c13-summit} \\
\bottomrule
\end{tabular}
\end{center}

\chaptersection{Related Topics}

\begin{itemize}[nosep, leftmargin=1.4em]
  \item \chref{c14} (Ch.~18, Fenwick Trees) --- backward/forward: Ch.~18 drew the
        expressiveness line (a BIT owns invertible prefix sums; arbitrary-range
        \(\min/\max\) under arbitrary writes is a segtree) and pointed it HERE;
        19.4's comparison box discharges that promise, and 19.3 contrasts the segtree
        descend against \(\coderef{c14-kth}\)'s lowbit-lift (no BIT re-teach).
  \item \chref{c15} (Ch.~20, Segment Trees II: Lazy) --- forward: the moment a whole
        interval is \emph{modified} (range add / range assign), the point-update tree
        of this chapter is not enough and lazy tags take over; 19.2 and 19.4 note
        exactly where that boundary falls but ship no lazy code.
  \item \chref{c19} (Ch.~21, Merge-Sort / Persistent Trees) --- forward: a \emph{pure}
        static \(\min/\max\) is \(O(1)\) via the sparse table, and ``\(k\)-th smallest
        in a static range'' needs the merge-sort or persistent tree; 19.5 forward-cites
        both rather than building either here.
  \item \chref{c16} (Ch.~22, Heavy-Light Decomposition) --- forward: path queries on a
        tree are Euler-flattened onto exactly this segment tree; mentioned by prose only.
  \item \chref{c24} (Ch.~5, Euler Tour) --- backward: owns \(\tin/\tout\) and the
        ``subtree \(=\) interval'' fact the 19.5 pointer reuses by label; the dynamic
        subtree engine is handed forward to Ch.~20/22.
  \item \chref{c11} (Ch.~11, Binary Lifting) --- backward: 19.3's descend is the
        tree-walk sibling of lifting by bits; back-cited, not re-derived.
  \item \chref{c01} (Ch.~1, complexity \& binary search) --- backward: the Summit's
        outer loop is the monotone binary search Ch.~1 frames; here the array being
        searched is the segment tree's answer axis.
  \item Ch.~6's treap section --- backward (prose only): it worked a range-sum with a
        self-balancing tree and forward-pointed to the segment-tree treatment 19.2
        supplies (LC 307 stays Ch.~6's link and is named, never re-linked).
\end{itemize}

\begin{lastword}[Last Word --- Chapter 19]
\emph{``Keep the whole interval tree, not just the lowbit shadows, and any
associative combine becomes a \(\bigO{\log n}\) walk: fold a range up from the
canonical cover, write a point down one root path, and -- when a stored aggregate can
answer `is the target under me?' -- descend root-to-leaf to find the very position a
Fenwick tree is blind to.''}

\medskip
The tree is one skeleton, many payloads. Store a scalar and you have range
\(\min/\max/\gcd\) under point set (\(\coderef{c13-query}\), the fold Ch.~18 could not
host). Store a four-field record and the max subarray folds in \(O(1)\) per merge, order
matters and the two iterative accumulators keep it straight (\(\thmref{c13-subarray}\),
\(\coderef{c13-iter-range}\)). Compact the recursion into a complete \(2n\) array and
the code shrinks to two loops and a half-open walk (\(\propref{c13-iter}\)). Read-only
and never updated, the same build is an oracle an offline sweep or a DP transition
calls a thousand times (\(\propref{c13-offline}\), CF~940E).

\medskip
Carry three reflexes forward:
\begin{enumerate}[nosep, leftmargin=1.6em]
  \item \emph{Fold vs descend.} A query that returns a \emph{value} is a fold over the
        cover; a query that returns an \emph{index} (first/leftmost/\(k\)-th) is a
        prune-as-you-go descent -- and only the tree can do the second
        (\(\thmref{c13-descend}\)).
  \item \emph{Identity is not optional.} Every fold needs the true neutral element;
        a disjoint branch returning anything but \(e\) is a silent wrong answer
        (\(\defref{c13-identity}\)).
  \item \emph{Composability is the payoff.} A segtree is the read engine inside a
        binary search (CF~514D, the Summit) and inside a DP (CF~940E); recognising the
        monotone predicate wrapped around a fold is the whole advanced game.
\end{enumerate}

\medskip
And three habits, drilled until boring: size the recursive array \(4n\) and the
complete array \(2n\) and know which you are typing; re-check the half-open/closed
convention and the descend's probe order against the smallest sample that could expose
them; and before writing a segment tree, ask whether your combine is a plain invertible
sum -- if it is, Chapter~18 already owns the cheaper answer.

\medskip
\textbf{Next:} \chref{c15} (Segment Trees II: Lazy Propagation) removes the last
restriction this chapter left standing -- instead of editing one leaf, you will add or
assign to a \emph{whole interval} and still answer range queries, by deferring writes
with lazy tags pushed down on demand. The point-update tree you can now build cold is
exactly the object that grows its shadow.
\end{lastword}
